\begin{helper}
Fix a candidate index consisting of a \(K\)-set \(I\subseteq \mathrm{Fin}(n)\)
and an \(S\)-tuple of base vertices, where
\[
K=\noderef{TwoBiteNaturalIndependenceScale}(1+\varepsilon,n).
\]
Let the covering set of base vertices be the image of this tuple, and split
that covering set into red centers \(B_R\) and blue centers \(B_B\).  For a
fixed embedding cell \(\pi\), let \(P_R\) and \(P_B\) be the red and blue
coordinate images of \(I\) under \(\pi\).  Consider the concrete candidate
event used in \(\noderef{MediumClosedPairsDeterministicWitnessCandidatePackage}\):
there is a witness set \(B\) contained in the covering set, with the stated
medium-class, size, and cutoff bounds, and whose normalized red/blue edge
count is larger than
\[
\frac{K n^{1/4-\varepsilon}}{6(\log n)^2}.
\]
If a sample \(\omega\) in this candidate has
\(\noderef{TwoBiteEmbedding}(\omega)=\pi\), then at least
\[
\left\lceil \frac{K n^{1/4-\varepsilon}}{6(\log n)^2}\right\rceil
\]
coordinates are present in the normalized coordinate set built from
\(B_R,B_B,P_R,P_B\), namely the coordinate set used by
\(\noderef{MediumClosedPairsCandidateCellTail}\).
\end{helper}
\begin{proof}
Unpack the candidate witness \(B\).  Its normalized-coordinate lower bound
says that the number of normalized red witness edges plus the number of
normalized blue witness edges is strictly larger than
\[
A=\frac{K n^{1/4-\varepsilon}}{6(\log n)^2}.
\]
By the defining property of the natural ceiling, this implies
\[
\lceil A\rceil
\le |\mathrm{RedEdges}(B)|+|\mathrm{BlueEdges}(B)|.
\]

It remains to compare these witness edge sets with the fixed coordinate set
formed from the covering tuple and the embedding cell.  Since the witness
set \(B\) is contained in the covering set, every red center from \(B\) lies
in \(B_R\), and every blue center from \(B\) lies in \(B_B\).  Since
\(\noderef{TwoBiteEmbedding}(\omega)=\pi\), the red and blue projections
\(\noderef{TwoBiteRedProjection}(\omega)\) and
\(\noderef{TwoBiteBlueProjection}(\omega)\) agree on \(I\) with the first and
second coordinates of \(\pi\), so the witness projection vertices lie in
\(P_R\) and \(P_B\).

Each ordered witness edge is an actual red or blue graph edge.  Hence its
endpoints are distinct, so after applying the same normalization map it
belongs to the corresponding diagonal-free red or blue raw coordinate set.
The definitions of \(\noderef{TwoBiteEdgeCoordinateValue}\),
\(\noderef{TwoBiteRedGraph}\), and \(\noderef{TwoBiteBlueGraph}\) show that
the normalized coordinate is present in \(\omega\); if normalization swaps
the two endpoints, symmetry of simple-graph adjacency supplies the same
edge in the reversed order.

Tagging red normalized edges by the left summand and blue normalized edges by
the right summand gives disjoint finite subsets of the present-coordinate
filter.  Their cardinalities are exactly the original red and blue normalized
edge counts.  Therefore the present-coordinate filter has cardinality at
least
\[
|\mathrm{RedEdges}(B)|+|\mathrm{BlueEdges}(B)|,
\]
and chaining this with the ceiling inequality proves the claimed threshold.
\end{proof}
